\chapter{Introduction}

\section{Contexte}

Les informations nécessaires à la compréhension d'un système ferroviaire ne sont pas réunies dans une seule source. Le métamodèle Arianes formalise les objets techniques et leurs relations, tandis que les fiches de procédure décrivent leur utilisation, les règles de gestion et les opérations à effectuer. Une recherche exclusivement textuelle retrouve des mots, mais maîtrise difficilement les dépendances entre objets. À l'inverse, un graphe représente ces dépendances sans contenir toutes les explications disponibles dans les documents.

Le projet vise à rapprocher ces deux formes de connaissances dans un assistant conversationnel. Celui-ci doit pouvoir répondre aussi bien à une question structurelle, telle que « quels sont les attributs de la classe Ligne ? », qu'à une question documentaire ou métier, telle que « quels équipements sont commandés par un poste d'aiguillage ? ».

\section{Problématique}

La problématique retenue est la suivante :

\begin{quote}
Comment exploiter conjointement le métamodèle Arianes et les fiches de procédure afin de produire des réponses techniques compréhensibles, traçables et limitées aux connaissances réellement disponibles ?
\end{quote}

Cette problématique impose plusieurs exigences :

\begin{itemize}
  \item représenter les relations du métamodèle dans Neo4j ;
  \item retrouver des passages même lorsque la question n'emploie pas exactement les mêmes mots ;
  \item distinguer une relation UML d'une connaissance fonctionnelle issue d'une fiche ;
  \item empêcher l'import automatique de relations métier incertaines ;
  \item conserver la source de chaque information ;
  \item isoler l'historique de chaque utilisateur et de chaque conversation.
\end{itemize}

\section{Objectifs}

Les objectifs opérationnels sont :

\begin{enumerate}
  \item convertir les sources XMI en données structurées ;
  \item construire un graphe technique du métamodèle ;
  \item découper et indexer les fiches documentaires ;
  \item mettre en place un retrieval GraphRAG déterministe ;
  \item développer une application web avec mémoire et authentification ;
  \item extraire une couche de concepts métier contrôlée ;
  \item évaluer la solution par des tests reproductibles.
\end{enumerate}

\section{Organisation du rapport}

Après les fondements théoriques, le rapport présente l'architecture générale, la construction du graphe technique, le fonctionnement du GraphRAG, puis le pipeline métier. Les derniers chapitres décrivent l'application, les résultats, les limites et les perspectives.

\attention{Les termes « extraction », « validation » et « import » désignent trois étapes distinctes. Une relation extraite par un LLM n'est pas considérée comme validée et n'est pas nécessairement présente dans Neo4j.}

